\scanpage{127}%
\mspage{121}%
% archive sheet number: 126
% Lines 1--16; the opening completes considération on scan 126.
ration préalable de $\Sigma$, le mal-aimé ? Par exemple, les supports \inserted{\unsure{constr.}} \uline{fermés} (\deleted{\unsure{associés}} \inserted{\unsure{relatifs}} à un magasin-\unsure{étalons|ateliers} $\unsure{\mathcal{F}|F}$) \unsure{peuvent-ils|pourront-ils} s'obtenir, en prenant l'ensemble ordonné déduit de $(\mathcal{F},\ll)$ en ``inversant'' les ``flèches'' $\unsure{\ll|\longleftarrow}$ (sur les subdivisions) ? Je ne suis pas $=$ sûr que la cat. des \unsure{fractions} obtenue soit bien un ens. ordonné, et encore moins (\unsure{même} s'il l'était) qu'il \deleted{\unsure{coïncide avec}} s'identifie à un sous-ens. ordonné de $\Sigma$. (Ce n'est pas là nécessairement un mal, si $\Sigma$ devait être \unsure{largué|légué} aux profits et pertes\ldots).

% Lines 17--23.
Je dirais quand-même, pour me \unsure{rassurer}, que dans les cas où $\Sigma$ est booléen (ateliers \unsure{ponctuaires}, et peut-être $=$ les ateliers fidèles ?, ou tout au moins \unsure{ainsi|aussi} les \unsure{maquettes}), il n'y a rien à lui reprocher. Ça inclut quand même déjà tous les cas qu'on peut
